%==============================================================================
%  Synapse — Documento técnico
%  El nodo como flecha: un modelo de tres coordenadas para el diseño de un
%  catálogo de nodos componible.
%
%  Compilar:  pdflatex Synapse_Node_Model.tex   (dos veces, para el índice)
%             o:  latexmk -pdf Synapse_Node_Model.tex
%==============================================================================
\documentclass[11pt,a4paper]{article}

%--- Codificación e idioma --------------------------------------------------
\usepackage[utf8]{inputenc}
\usepackage[T1]{fontenc}
\usepackage[spanish,es-noquoting]{babel}
\usepackage{lmodern}
\usepackage{microtype}

%--- Geometría de página ----------------------------------------------------
\usepackage[a4paper,top=2.6cm,bottom=2.6cm,left=2.8cm,right=2.8cm]{geometry}
\usepackage{setspace}
\onehalfspacing

%--- Utilidades -------------------------------------------------------------
\usepackage{amsmath,amssymb}
\usepackage{amsthm}
\usepackage{graphicx}
\usepackage{booktabs}
\usepackage{enumitem}
\usepackage{abstract}
\usepackage{titlesec}
\usepackage{fancyhdr}
\usepackage{tikz}
\usetikzlibrary{arrows.meta,calc,positioning,fit,backgrounds,shapes.geometric}
\usepackage{pifont}           % \ding (marcas de verificación)
\usepackage[hidelinks]{hyperref}
\usepackage{xcolor}

%--- Acento discreto de marca -----------------------------------------------
\definecolor{synapseCyan}{HTML}{0E7490}
\definecolor{synapseInk}{HTML}{212121}
\definecolor{synapseSoft}{HTML}{E4F1F4}

%--- Estilo de títulos de sección -------------------------------------------
\titleformat{\section}
  {\normalfont\Large\bfseries\color{synapseInk}}{\thesection}{0.6em}{}
\titleformat{\subsection}
  {\normalfont\large\bfseries\color{synapseInk}}{\thesubsection}{0.6em}{}

%--- Encabezado / pie -------------------------------------------------------
\pagestyle{fancy}
\fancyhf{}
\renewcommand{\headrulewidth}{0.4pt}
\fancyhead[L]{\small\itshape Synapse — El nodo como flecha}
\fancyhead[R]{\small\thepage}

%--- Metadatos del PDF ------------------------------------------------------
\hypersetup{
  pdftitle={El nodo como flecha: un modelo de tres coordenadas},
  pdfauthor={Mauricio Daniel Klainbard},
  pdfsubject={Diseno de un catalogo de nodos componible},
  colorlinks=false
}

%--- Entornos ---------------------------------------------------------------
\theoremstyle{definition}
\newtheorem{definicion}{Definición}[section]
\newtheorem{ejemplo}{Ejemplo}[section]
\theoremstyle{plain}
\newtheorem{proposicion}{Proposición}[section]
\newtheorem{observacion}{Observación}[section]

\newcommand{\prop}[1]{\medskip\noindent\textbf{Propiedad.}\enspace\emph{#1}\medskip}
\newcommand{\mundo}[1]{\textsc{#1}}   % nombre de un mundo de datos

%--- Estilos TikZ -----------------------------------------------------------
\tikzset{
  box/.style   ={draw, rounded corners, align=center, font=\small, inner sep=5pt},
  world/.style ={draw, ellipse, align=center, font=\small,
                 minimum width=2.3cm, minimum height=1.0cm},
  hub/.style   ={draw, ellipse, very thick, align=center, font=\small,
                 minimum width=2.7cm, minimum height=1.3cm, fill=synapseSoft},
  leaf/.style  ={draw, ellipse, dashed, align=center, font=\small,
                 minimum width=2.3cm, minimum height=1.0cm},
  dot/.style   ={circle, draw, fill=synapseInk, inner sep=1.6pt},
  ar/.style    ={-{Stealth[length=2.2mm]}, thick},
  no/.style    ={dashed, thick, gray},
}

\begin{document}

%============================================================================
%  CARÁTULA
%============================================================================
\begin{titlepage}
\centering
\vspace*{1.5cm}

{\color{synapseCyan}\rule{\linewidth}{1.2pt}}\\[0.5cm]

{\Huge\bfseries El nodo como flecha\par}
\vspace{0.4cm}
{\LARGE Un modelo de tres coordenadas para el diseño\\ de un catálogo de nodos componible\par}
\vspace{0.3cm}
{\large\itshape Familias, cardinalidad y mundos de datos\par}

\vspace{0.4cm}
{\color{synapseCyan}\rule{\linewidth}{1.2pt}}\\[1.6cm]

{\large Documento técnico\par}
\vspace{0.2cm}
{\normalsize Motor de automatización de flujos de trabajo\par}

\vfill

{\large Mauricio Daniel Klainbard\par}
\vspace{0.3cm}
{\normalsize\today\par}
\vspace{1.2cm}

{\footnotesize
Documento de carácter metodológico. Presenta un modelo conceptual para el diseño
de nodos y su clasificación; es una arquitectura de razonamiento, no una
especificación de implementación, y no altera los nodos del catálogo actual.\par}

\end{titlepage}

%============================================================================
%  RESUMEN
%============================================================================
\thispagestyle{fancy}
\begin{abstract}
\noindent
Un catálogo de nodos de automatización se diseña, habitualmente, nodo por nodo y
de forma aislada. Este documento sostiene que el diseño aislado es la causa raíz
de una patología recurrente: nodos que, al integrarse al ecosistema, exhiben pocas
interacciones útiles con el resto y obligan a rediseños sucesivos. Como remedio se
propone un \emph{modelo de diseño} que reemplaza la pregunta ``¿qué hace este
nodo?'' por ``¿cómo se compone este nodo con los demás?''. El punto de partida es
un intento fallido pero fructífero: definir las \emph{familias} de nodos como
clases de equivalencia de una relación de no-conexión. Se muestra que dicha
relación es reflexiva y simétrica pero \emph{no transitiva}, por lo que no induce
una partición. El análisis de por qué falla conduce al modelo central: un nodo no
es un punto sino una \emph{flecha} entre \emph{mundos de datos}, y su lugar en el
ecosistema queda fijado por tres coordenadas ---la firma de mundo, que fija a la
vez su origen y su destino; la cardinalidad; y la aridad de puertos que describe
su ramificación---. Se demuestra que este esquema recupera la intuición original de
familia como un teorema en lugar de una definición, clasifica correctamente los
nodos de control de flujo (\texttt{if}, \texttt{switch}, \texttt{merge}) y, sobre
todo, ofrece una disciplina de diseño operativa. Se argumenta finalmente que las
tres coordenadas constituyen una \emph{gramática composicional} que eleva de forma
notable la calidad de los flujos que un asistente de lenguaje ---incluso uno
pequeño--- puede construir.

\vspace{0.4cm}
\noindent\textbf{Palabras clave:} diseño de nodos, composicionalidad, relaciones
de equivalencia, cardinalidad, tipos de datos, automatización de flujos de trabajo,
asistentes de lenguaje.
\end{abstract}

\vspace{0.5cm}
\tableofcontents
\newpage

%============================================================================
\section{La dificultad de diseñar un nodo}
%============================================================================

\subsection{El diseño aislado y su patología}

Diseñar un nodo parece, a primera vista, un problema local: se elige una operación
útil ---renderizar un gráfico, resolver una ecuación diferencial, leer una hoja de
cálculo---, se definen sus parámetros de configuración y su lógica, y se lo agrega
al catálogo. Esta forma de proceder, que llamaremos \emph{diseño aislado}, trata
cada nodo como una unidad autocontenida cuya bondad se juzga por lo bien que
realiza su tarea.

El diseño aislado falla por una razón estructural: un nodo no vive solo. Su valor
en una plataforma de automatización no reside en lo que hace por sí mismo, sino en
\emph{con qué otros nodos se compone}. Un nodo excelente en aislamiento pero que no
recibe datos de ninguna fuente natural, o cuya salida ningún otro nodo puede
consumir, es un nodo pobre: una isla en el grafo. Y esta propiedad ---la riqueza de
sus interacciones--- es justamente la que el diseño aislado no contempla, porque
mira el nodo y no el ecosistema.

La consecuencia práctica es un ciclo de retrabajo. Un nodo concebido de forma
aislada se implementa, se intenta insertar en flujos reales, y solo entonces se
descubre que su forma de entrada no coincide con lo que producen los nodos que
deberían precederlo, o que su salida obliga a intercalar adaptadores, o que su
cardinalidad hace estallar la tubería. Se lo rediseña; el rediseño revela una nueva
fricción; se lo vuelve a rediseñar. El costo no está en escribir el nodo sino en
\emph{descubrir, tarde y empíricamente, cómo debería haber sido su forma}.

\subsection{Una metodología insuficiente: un nodo por día}

Una respuesta natural al problema es organizativa: fijar un ritmo ---por ejemplo, un
nodo por día--- y avanzar. Pero el ritmo no ataca la causa. Diseñar un nodo por día
\emph{sigue siendo} diseño aislado, solo que calendarizado; cada nodo se concibe
sin un marco que lo obligue a declarar, de entrada, cómo encaja con los demás. El
retrabajo no desaparece: se difiere y se acumula.

Lo que hace falta no es una cadencia sino un \emph{modelo}: un lenguaje pequeño y
preciso en el que la pregunta ``¿cómo se conecta este nodo con el resto del
ecosistema?'' se pueda responder \emph{antes} de escribir una sola línea, y en el
que un mal diseño ---una isla--- sea visible como tal en el papel. El resto de este
documento construye ese modelo. Es deliberadamente una arquitectura de
razonamiento, externa al código: no introduce propiedades nuevas en los nodos ni
mecanismos de verificación en el motor, y no altera en nada el catálogo actual. Su
único propósito es enseñarnos a diseñar.

%============================================================================
\section{Un primer intento: las familias como clases de equivalencia}
%============================================================================

\subsection{La intuición de familia}

La primera observación fructífera es que los nodos se agrupan en \emph{familias}
definidas por el ecosistema sobre el que operan: la familia de los gráficos
(\mundo{Chart}), la de las integraciones con un proveedor (\mundo{Google}), la de
los métodos numéricos. Se buscó dar a esta noción intuitiva una base formal. La
herramienta evidente, tomada de la teoría de conjuntos, es la \emph{relación de
equivalencia}: una relación que sea reflexiva, simétrica y transitiva induce una
partición del conjunto en clases disjuntas~\cite{hrbacek}, y esas clases serían,
precisamente, las familias.

\subsection{La relación de no-conexión}

Como diseñábamos pensando en la interacción entre familias, se propuso definir la
relación por su \emph{ausencia}. La observación de partida fue que dos nodos de la
familia \mundo{Chart} no pueden conectarse entre sí: todos producen una imagen a
partir de datos numéricos, y ninguno consume una imagen ni produce datos numéricos,
de modo que la salida de uno nunca alimenta la entrada de otro.

\begin{definicion}[Relación de no-conexión]
Sean $a,b$ dos nodos. Se dice que $a \perp b$ si y solo si no existe conexión
posible entre $a$ y $b$ en ninguna de las dos direcciones; es decir, ni la salida
de $a$ puede alimentar significativamente la entrada de $b$, ni la de $b$ la de
$a$.
\end{definicion}

La idea era sembrar cada familia con un nodo representativo $r$ y definir su clase
como $\{\,x : x \perp r\,\}$. Examinemos si $\perp$ es una relación de equivalencia.

\paragraph{Simetría.} Es inmediata por construcción. La cláusula ``en ninguna de
las dos direcciones'' es indiferente al orden de $a$ y $b$; luego
$a \perp b \iff b \perp a$.

\paragraph{Reflexividad.} A primera vista falla. Un nodo \texttt{filter} puede
conectarse con otro \texttt{filter} ---la salida del primero es una entrada válida
para el segundo---, de modo que $\lnot(\texttt{filter} \perp \texttt{filter})$. Sin
embargo, dos nodos de la misma especie encadenados son \emph{reducibles} a uno
solo: dos filtros en serie equivalen a un único filtro con la conjunción de ambas
condiciones. Bajo la convención de que dos instancias reducibles se identifican, la
autoconexión no aporta un grafo genuinamente nuevo, y podemos adoptar
$a \perp a$ por convención. La reflexividad se sostiene, aunque como estipulación y
no como hecho.

\subsection{El fracaso: la transitividad}

La propiedad decisiva ---y la que hunde el modelo--- es la transitividad. Para que
$\perp$ particione, debería cumplirse
\[
  (a \perp b) \;\wedge\; (b \perp c) \;\Longrightarrow\; (a \perp c).
\]
No se cumple. La razón profunda es que ``no conectar'' es la \emph{negación} de una
relación de compatibilidad que es a la vez \emph{dirigida} (la salida de uno contra
la entrada del otro) y de \emph{dos ejes} (qué produce cada nodo, qué consume cada
nodo); la negación de una relación con esa estructura no es transitiva en general.
Un contraejemplo mínimo lo exhibe.

\begin{ejemplo}[Ruptura de la transitividad]\label{ej:trans}
Considérense tres nodos:
\begin{itemize}[nosep,leftmargin=1.4em]
  \item $A$: una \emph{fuente numérica} (produce números, no consume nada);
  \item $B$: una \emph{fuente de imágenes} (produce una imagen, no consume nada);
  \item $C$: un \texttt{chart} (consume números, produce una imagen).
\end{itemize}
Entonces:
\begin{itemize}[nosep,leftmargin=1.4em]
  \item $A \perp B$: dos fuentes, ninguna consume la salida de la otra. \checkmark
  \item $B \perp C$: $C$ consume números y $B$ produce una imagen (incompatible);
        $B$ no consume nada. \checkmark
  \item $\lnot(A \perp C)$: $A$ produce números y $C$ los consume, luego
        $A \to C$ es una conexión legítima. \ding{55}
\end{itemize}
Se tiene $A \perp B$ y $B \perp C$ pero $\lnot(A \perp C)$: la transitividad falla.
\end{ejemplo}

\begin{figure}[h]
\centering
\begin{tikzpicture}
  \node[box] (A) at (0,0)   {$A$\\[2pt]\footnotesize $\varnothing \to$ número};
  \node[box] (B) at (6,0)   {$B$\\[2pt]\footnotesize $\varnothing \to$ imagen};
  \node[box] (C) at (3,-2.6){$C$: \texttt{chart}\\[2pt]\footnotesize número $\to$ imagen};
  \draw[no] (A) -- node[above]{$\perp$} (B);
  \draw[no] (B) -- node[above right]{$\perp$} (C);
  \draw[ar,synapseCyan] (A) -- node[above left]{conecta} (C);
\end{tikzpicture}
\caption{El contraejemplo del Ejemplo~\ref{ej:trans}. Las aristas punteadas indican
no-conexión ($\perp$); la flecha llena, una conexión legítima. $A\perp B$ y
$B\perp C$ no fuerzan $A\perp C$.}
\label{fig:trans}
\end{figure}

Hay, además, una segunda grieta que revela que $\perp$ ni siquiera captura la
noción que pretendía. La familia \mundo{Chart} es, en efecto, cerrada bajo $\perp$
(sus miembros no se conectan entre sí), pero la familia \mundo{Google} no lo es:
leer un correo y escribir en una hoja de cálculo son ambos nodos de \mundo{Google}
y \emph{sí} se conectan. La relación $\perp$ describe un \emph{tipo} de familia ---la
conversora--- y no otro ---la de proveedor---. La misma causa (la compresión de una
estructura dirigida y de dos ejes en una relación simétrica de un solo eje) produce
las dos fallas.

\begin{observacion}
El intento no fue estéril. Exigirle a la noción de familia las tres propiedades de
una relación de equivalencia ---la disciplina que Hrbacek y Jech~\cite{hrbacek}
imponen a toda partición--- es exactamente lo que expuso que $\perp$ era el objeto
equivocado. Y la teoría de conjuntos ofrece también el remedio: toda función induce
una relación de equivalencia canónica ---su \emph{núcleo}, $a \sim b \iff f(a)=f(b)$
---, cuyas clases son las fibras de la función y particionan sin excepción. El
modelo que sigue no abandona la idea de equivalencia: la funda sobre la función
correcta.
\end{observacion}

%============================================================================
\section{El modelo de tres coordenadas: el nodo como flecha}
%============================================================================

\subsection{El nodo como flecha entre mundos de datos}

El error de fondo del primer intento fue tratar el nodo como un \emph{punto} que se
relaciona con otros puntos. La corrección es tratarlo como una \emph{flecha}: un
nodo transporta datos de un \emph{mundo} a otro.

\begin{definicion}[Nodo como flecha]
Sea $\mathcal{W}$ el conjunto de \emph{mundos de datos} (la forma semántica que
puede tener aquello que fluye: una tabla de registros, un texto, una imagen, un
archivo, la nada, etc.; la Sección~5 los desarrolla). Un nodo $f$ es una flecha
\[
  f : X \longrightarrow Y, \qquad X,Y \in \mathcal{W},
\]
donde $X=\operatorname{src}(f)$ es su mundo de origen (lo que consume) y
$Y=\operatorname{dst}(f)$ su mundo de destino (lo que produce).
\end{definicion}

\begin{figure}[h]
\centering
\begin{tikzpicture}
  \node (x) at (0,0) {$X$};
  \node (y) at (5,0) {$Y$};
  \draw[ar] (x) -- node[above]{$f$} (y);
  \node[font=\footnotesize] at (2.5,-0.55) {origen $\to$ destino};
\end{tikzpicture}
\caption{La unidad del modelo: un nodo es una flecha de un mundo de origen a un
mundo de destino.}
\label{fig:arrow}
\end{figure}

Bajo esta lente, la noción de conexión se vuelve precisa y recupera lo que $\perp$
buscaba a tientas.

\begin{definicion}[Conexión componible]
Dos nodos $a,b$ se conectan, $a \to b$, de forma significativa si y solo si
$\operatorname{dst}(a) = \operatorname{src}(b)$: la salida de $a$ vive en el mismo
mundo que la entrada de $b$. La conexión es, literalmente, la \emph{composición de
flechas}.
\end{definicion}

Esto es fiel a la naturaleza del motor, cuyo cable es \emph{no tipado}: por él viaja
siempre un elemento con una forma libre, y el editor jamás prohíbe una arista. La
compatibilidad no la impone el grafo; la impone la \emph{forma semántica} ---el
mundo--- de lo que la flecha produce frente a lo que la siguiente consume. El modelo
no describe una restricción del motor, sino la disciplina con la que \emph{nosotros}
razonamos sobre él.

\subsection{La primera coordenada y su insuficiencia}

Con las flechas aparece de inmediato la relación de equivalencia correcta, como
núcleo de una función.

\begin{definicion}[Firma de mundo]
La \emph{firma de mundo} de un nodo es el par
$\sigma(f) = (\operatorname{src}(f), \operatorname{dst}(f))$. Dos nodos son de la
misma familia-de-mundo cuando $\sigma(a)=\sigma(b)$.
\end{definicion}

Por ser el núcleo de $\sigma$, esta relación es reflexiva, simétrica y transitiva
\emph{por construcción}: particiona sin colapsar, y el escenario de colapso que
temíamos ---un nodo que arrastra dos familias a fusionarse--- es imposible, porque
las fibras de una función nunca se solapan. Además, la propiedad que $\perp$ quería
\emph{definir} reaparece aquí como un \emph{teorema}: una familia con
$\operatorname{src}\neq\operatorname{dst}$ (fuera de la diagonal) es automáticamente
cerrada bajo no-conexión ---dos flechas $X\to Y$ con $X\neq Y$ no componen, pues
requeriría $Y=X$---. La familia \mundo{Chart}, toda ella $\mundo{Tabla}\to\mundo{Imagen}$,
lo satisface. La familia \mundo{Google}, que abarca varias firmas de mundo, no; y
esa es exactamente la razón por la que sus nodos sí interactúan.

Pero una sola coordenada es \emph{insuficiente}, y el propio catálogo lo delata.
Considérense tres operaciones sobre tablas:
\[
  \texttt{sort} : \mundo{Tabla}\to\mundo{Tabla},\quad
  \texttt{filter} : \mundo{Tabla}\to\mundo{Tabla},\quad
  \texttt{reduce} : \mundo{Tabla}\to\mundo{Tabla}.
\]
Las tres comparten firma de mundo: caen en la misma clase, la \emph{diagonal}
$\mundo{Tabla}\to\mundo{Tabla}$. La firma de mundo sugiere entonces que son la misma
familia y ---peor--- que podrían reducirse a un solo nodo, exactamente como dos
filtros en serie se reducían a uno. Y sin embargo no lo son: ordenar, filtrar y
plegar son operaciones distintas que no se sustituyen entre sí. Falta una
coordenada que las separe.

\subsection{La segunda coordenada: la cardinalidad}

Lo que distingue a \texttt{sort}, \texttt{filter} y \texttt{reduce} no es a dónde
llevan el dato, sino \emph{cómo transforman la cantidad de elementos}. Esta es la
segunda coordenada, y es de una clase especialmente sólida: no es vocabulario
elegido, sino una invariante estructural que se lee de la operación misma.

\begin{definicion}[Cardinalidad]
La \emph{cardinalidad} $\kappa(f)$ de un nodo describe la relación entre el número
$N$ de elementos que entran y el número $M$ que salen. Las clases relevantes son:
\begin{itemize}[nosep,leftmargin=1.4em]
  \item \textbf{preserva} ($N{:}N$): un elemento entra, un elemento sale
        (\texttt{sort}, \texttt{set});
  \item \textbf{contrae selectivo} ($N{:}M,\ M\le N$): descarta elementos
        (\texttt{filter});
  \item \textbf{contrae total} ($N{:}1$): pliega el lote en un resultado
        (\texttt{reduce});
  \item \textbf{expande} ($N{:}M,\ M\ge N$): un elemento genera varios
        (\texttt{split\_out});
  \item \textbf{fuente} ($0{:}N$) y \textbf{sumidero} ($N{:}0$): nace o muere el
        stream.
\end{itemize}
\end{definicion}

\begin{table}[h]
\centering
\small
\begin{tabular}{@{}llll@{}}
\toprule
Nodo & Firma de mundo & Cardinalidad & Rol \\
\midrule
\texttt{sort}       & \mundo{Tabla}$\to$\mundo{Tabla} & $N{:}N$          & preserva (permuta) \\
\texttt{set}        & \mundo{Tabla}$\to$\mundo{Tabla} & $N{:}N$          & preserva (transforma por ítem) \\
\texttt{filter}     & \mundo{Tabla}$\to$\mundo{Tabla} & $N{:}M,\ M\le N$ & contrae selectivo \\
\texttt{split\_out} & \mundo{Tabla}$\to$\mundo{Tabla} & $N{:}M,\ M\ge N$ & expande \\
\texttt{reduce}     & \mundo{Tabla}$\to$\mundo{Tabla} & $N{:}1$          & contrae total \\
\bottomrule
\end{tabular}
\caption{La cardinalidad separa la diagonal $\mundo{Tabla}\to\mundo{Tabla}$ que la
firma de mundo confundía. (Persiste un residuo: \texttt{sort} y \texttt{set} siguen
compartiendo $N{:}N$; los separa una distinción ulterior, la de transformar el orden
frente al contenido. La cardinalidad es necesaria, no siempre suficiente: es una
coordenada más, no la última.)}
\label{tab:card}
\end{table}

La cardinalidad no solo refina la clasificación: \emph{compone}, y es ahí donde se
vuelve una herramienta de diseño y no una mera etiqueta. Si se piensa la
cardinalidad como un multiplicador sobre el tamaño del lote ---contrae ($<1$),
preserva ($=1$), expande ($>1$)---, al encadenar nodos los multiplicadores se
multiplican:
\[
  \text{preserva}\circ X = X, \qquad
  \text{expande}\circ\text{expande} = \text{expande}, \qquad
  \text{expande}\ \text{seguida de}\ (N{:}1) = 1.
\]
De aquí se leen tres consecuencias de diseño. Primera: dos expansores en serie son
la firma de una \emph{explosión} combinatoria del lote; el modelo la anticipa en el
papel. Segunda: la reflexividad-por-reducibilidad del intento inicial se precisa
---se sostiene \emph{dentro de una misma clase de cardinalidad} ($\texttt{sort}\circ
\texttt{sort}=\texttt{sort}$, $\texttt{filter}\circ\texttt{filter}=\texttt{filter}$),
pero $\texttt{reduce}\circ\texttt{reduce}$ ya es degenerado---. Tercera: aparece un
\emph{arco de tubería} canónico, \emph{expandir $\to$ preservar$^*$ $\to$ contraer},
que es la forma de todo pipeline sano (Figura~\ref{fig:stratum}).

\begin{figure}[h]
\centering
\begin{tikzpicture}[node distance=6mm]
  \node[box, fill=synapseSoft] (s)  {fuente\\\footnotesize $0{:}N$};
  \node[box, right=1.4cm of s]  (p)  {preservadores\\\footnotesize $N{:}N$};
  \node[box, right=1.4cm of p]  (c)  {contractor\\\footnotesize $N{:}1$};
  \node[box, right=1.4cm of c, fill=synapseSoft]  (k)  {sumidero\\\footnotesize $N{:}0$};
  \draw[ar] (s) -- (p);
  \draw[ar] (p) -- (c);
  \draw[ar] (c) -- (k);
\end{tikzpicture}
\caption{El estrato de tubería que induce la cardinalidad, ortogonal al eje de
mundos: las fuentes abren trabajo, los preservadores lo transforman sin cambiar el
conteo, los contractores lo pliegan, los sumideros lo consumen.}
\label{fig:stratum}
\end{figure}

%============================================================================
\section{El modelo frente a \texttt{if} y \texttt{switch}}
%============================================================================

Los nodos de control de flujo son la prueba de fuego del modelo, porque no son una
sola flecha: \texttt{if} tiene dos puertos de salida y \texttt{switch} tiene puertos
dinámicos. Lejos de romper el esquema, obligan a explicitar la tercera coordenada
---la aridad de puertos--- que el resto de los nodos usaba de forma trivial.

\subsection{En el eje de mundos: son la identidad}

Lo primero es que \texttt{if} y \texttt{switch} \emph{no tocan el elemento}: entra
una tabla, sale la misma tabla, sin transformar, por cada puerto. Son enrutadores
puros. En la firma de mundo son endomorfismos identidad
$\mundo{Tabla}\to\mundo{Tabla}$, como \texttt{sort}. No mueven el dato a otro mundo;
lo mueven a otro \emph{cable}.

\subsection{En el eje de cardinalidad: una partición conservativa}

Como caja negra son $N{:}N$: no crean ni destruyen elementos. Pero ese $N$ se
\emph{reparte} entre $k$ puertos de salida,
\[
  N = n_1 + n_2 + \dots + n_k, \qquad n_i \ge 0,
\]
y cada puerto, mirado por separado, es $N{:}n_i$ con $n_i \le N$: es decir, \emph{un
\texttt{filter}}. De aquí la unificación que confirma que el modelo es el correcto:
\begin{quote}
\texttt{switch} es un banco de \texttt{filter} en paralelo que particiona el lote;
\texttt{if} son dos \texttt{filter} complementarios; \texttt{filter} es un
\texttt{if} al que se le descartó el puerto negativo.
\end{quote}
\texttt{filter}, \texttt{if} y \texttt{switch} son la \emph{misma operación} a
distinta aridad de puertos ($1$ que se conserva, $2$, $k$). Lo único que cambia es
qué se hace con el resultado del predicado: descartarlo, rutearlo a dos puertos, o a
$k$ (con un puerto por defecto que hace de la partición una función total sobre el
elemento).

\begin{figure}[h]
\centering
\begin{tikzpicture}[>=Stealth]
  % if
  \begin{scope}
    \node[dot] (ii) at (0,0) {};
    \node[right=1.6cm of ii] (it) at (1.6,0.7) {};
    \node (if) at (1.6,-0.7) {};
    \draw[ar] (ii) -- (1.6,0.7) node[right,font=\footnotesize]{true};
    \draw[ar] (ii) -- (1.6,-0.7) node[right,font=\footnotesize]{false};
    \node[font=\small\bfseries] at (0.5,-1.6) {\texttt{if} (1$\to$2)};
  \end{scope}
  % switch
  \begin{scope}[xshift=5cm]
    \node[dot] (si) at (0,0) {};
    \draw[ar] (si) -- (1.6,1.0) node[right,font=\footnotesize]{ruta $1$};
    \draw[ar] (si) -- (1.6,0.3) node[right,font=\footnotesize]{ruta $2$};
    \draw[ar] (si) -- (1.6,-0.4) node[right,font=\footnotesize]{$\vdots$};
    \draw[ar] (si) -- (1.6,-1.1) node[right,font=\footnotesize]{defecto};
    \node[font=\small\bfseries] at (0.5,-1.9) {\texttt{switch} (1$\to k$)};
  \end{scope}
  % merge
  \begin{scope}[xshift=10.5cm]
    \node[dot] (mo) at (1.8,0) {};
    \draw[ar] (0,1.0) node[left,font=\footnotesize]{$X_1$} -- (mo);
    \draw[ar] (0,0.3) node[left,font=\footnotesize]{$X_2$} -- (mo);
    \draw[ar] (0,-0.4) node[left,font=\footnotesize]{$\vdots$} -- (mo);
    \draw[ar] (0,-1.1) node[left,font=\footnotesize]{$X_k$} -- (mo);
    \node[font=\small\bfseries] at (1.3,-1.9) {\texttt{merge} ($k\to$1)};
  \end{scope}
\end{tikzpicture}
\caption{La aridad de puertos como tercera coordenada estructural. \texttt{if} y
\texttt{switch} ramifican (una entrada, $k$ salidas, partición conservativa);
\texttt{merge} es su dual exacto (fusiona $k$ entradas en una salida, por unión).}
\label{fig:ports}
\end{figure}

\subsection{La coordenada que esto fuerza, y su dual}

Si solo tuviéramos las dos coordenadas anteriores, \texttt{if}/\texttt{switch}
colapsarían con \texttt{sort} (todos $N{:}N$, $\mundo{Tabla}\to\mundo{Tabla}$). Lo
que los separa es la \emph{aridad de puertos}: single-port frente a multi-port. Y
esta coordenada tiene un dual exacto que cierra el modelo con elegancia:
\texttt{merge} es el co-enrutador ---$k$ puertos de entrada $\to$ uno de salida,
unión en lugar de partición---. \texttt{switch} y \texttt{merge} son las dos caras
de una misma moneda: abrir el flujo por predicado y cerrarlo por estrategia.

La familia de control de flujo queda así definida con una firma nítida y
principista: \emph{identidad en el mundo}, \emph{multi-puerto}, y \emph{cardinalidad
conservativa} (partición al ramificar, unión al fusionar). Son los nodos que
reconfiguran la \emph{topología} del flujo sin tocar jamás el dato ---y eso es lo que
los distingue de la diagonal de transformación, que es single-port---. Los puertos
dinámicos de \texttt{switch} no son una excepción: su aridad $k$ es un parámetro de
configuración, el mismo patrón ``configuración $\to$ forma'' que ya gobierna la
cardinalidad de otros nodos.

Se observa el patrón que estructura todo el modelo: \emph{cada coordenada resuelve
un colapso que la anterior no podía}. La firma de mundo separó \mundo{Chart} de
\texttt{filter}, pero fundió \texttt{sort}/\texttt{filter}/\texttt{reduce}; la
cardinalidad los separó, pero fundió \texttt{sort} con \texttt{if}; la aridad de
puertos separa el enrutador del transformador. \texttt{if} y \texttt{switch} no son
una excepción al modelo: son lo que lo completa.

%============================================================================
\section{Los mundos de datos y la disciplina de diseño}
%============================================================================

\subsection{Qué es un mundo}

La primera coordenada descansa sobre el conjunto $\mathcal{W}$ de mundos de datos.
Un mundo es la \emph{forma semántica} de lo que fluye, no su serialización ni su
tipo en el motor. Es un vocabulario que vive en la cabeza del diseñador; no se
declara en ningún archivo ni altera nodo alguno. Un catálogo de trabajo usa unos
pocos:

\begin{itemize}[leftmargin=1.4em]
  \item \mundo{Tabla} --- una colección de elementos con campos. Es el mundo central,
        el gran concentrador: casi todo produce o consume tablas.
  \item \mundo{Texto} --- una cadena sin estructura impuesta; aquí viven el análisis
        de CSV/JSON y las expresiones regulares.
  \item \mundo{Archivo} --- bytes opacos referenciados (un adjunto); mundo de
        entrada de los nodos de documentos y visión.
  \item \mundo{Imagen} y \mundo{Documento} --- artefactos renderizados (un PNG, un
        PDF); mundos casi terminales.
  \item \mundo{Escalar} --- un valor único, el resultado de un pliegue.
  \item \mundo{Recurso} --- un asa a un objeto de un proveedor (una fila de una hoja
        de Google, un mensaje de Slack).
  \item \mundo{Vacío} ($\varnothing$) --- la nada: el origen de las fuentes ($0{:}N$)
        y el destino de los sumideros ($N{:}0$).
\end{itemize}

\begin{figure}[h]
\centering
\begin{tikzpicture}[>=Stealth]
  \node[hub]   (tab) at (0,0)      {\mundo{Tabla}};
  \node[world] (txt) at (-4.6,1.7) {\mundo{Texto}};
  \node[world] (vac) at (-4.6,-1.7){$\varnothing$};
  \node[world] (esp) at (0,3.0)    {\mundo{Especif.}};
  \node[leaf]  (img) at (4.6,1.7)  {\mundo{Imagen}};
  \node[leaf]  (arc) at (4.6,-1.7) {\mundo{Archivo}};
  \node[world] (esc) at (0,-3.0)   {\mundo{Escalar}};

  \draw[ar] (txt) to[bend left=8]  node[above,font=\scriptsize]{\texttt{csv\_parse}} (tab);
  \draw[ar] (tab) to[bend left=8]  node[below,font=\scriptsize]{\texttt{csv\_build}} (txt);
  \draw[ar] (vac) -- node[below,sloped,font=\scriptsize]{fuentes} (tab);
  \draw[ar] (esp) -- node[right,font=\scriptsize]{\texttt{ode\_solve}} (tab);
  \draw[ar] (tab) -- node[above,sloped,font=\scriptsize]{\texttt{chart}} (img);
  \draw[ar] (tab) -- node[below,sloped,font=\scriptsize]{\texttt{excel\_write}} (arc);
  \draw[ar] (tab) -- node[right,font=\scriptsize]{\texttt{reduce}} (esc);
  \draw[ar,synapseCyan] (tab) to[out=-35,in=35,looseness=6]
        node[right,font=\scriptsize]{\texttt{filter, sort, set, if, switch}} (tab);
\end{tikzpicture}
\caption{El mapa de mundos. Los nodos son las flechas. \mundo{Tabla} es el
concentrador (muchas entradas y salidas); \mundo{Imagen} y \mundo{Archivo} son hojas
(solo reciben); el lazo sobre \mundo{Tabla} reúne la diagonal de transformación y
control. Diseñar un nodo es elegir dónde poner su flecha en este mapa.}
\label{fig:worlds}
\end{figure}

\subsection{Por qué los mundos guían el diseño}

El mapa de mundos convierte ``diseñar un nodo rico en interacciones'' en una
instrucción operativa. Un nodo es rico cuando su flecha cae en un \emph{cruce de
alto tráfico}: su origen es un mundo que muchas familias producen (un destino
popular) y su destino es un mundo que muchas familias consumen (un origen popular).
\mundo{Tabla} es el cruce por excelencia; un nodo que atraviesa \mundo{Tabla} en
cualquiera de sus dos puntas hereda, sin esfuerzo, la posibilidad de componerse con
todo el resto del catálogo.

El mapa también hace visible el \emph{antipatrón} de un solo vistazo. Un nodo cuya
flecha va de un \mundo{Mundo\ exótico} a otro \mundo{Mundo\ exótico} ---un mundo que
ningún otro nodo produce, hacia un mundo que ningún otro consume--- es una isla:
por bien que realice su tarea, nada puede alimentarlo y nada puede leer su salida.
Cero interacciones posibles.

\prop{Un nodo $f:X\to Y$ tiene interacciones potenciales no vacías si y solo si
existe algún nodo cuyo destino sea $X$ (un proveedor de su entrada) y algún nodo
cuyo origen sea $Y$ (un consumidor de su salida). El diseño aislado no puede
verificar esta condición porque no mira el resto del mapa; el modelo de mundos la
vuelve una comprobación previa a escribir el nodo.}

La regla de diseño se enuncia entonces sin ambigüedad: \emph{al menos una de las dos
puntas de la flecha debe apoyarse en un mundo con tráfico}. Si ambas puntas son
exóticas, el nodo no debe diseñarse todavía; primero hay que poblar el mundo
exótico con nodos que lo produzcan o lo consuman, o bien reformular la operación
para que una punta aterrice en un concentrador. Esta es, exactamente, la deliberación
que el diseño aislado omitía y que causaba el ciclo de retrabajo de la Sección~1.

%============================================================================
\section{Consecuencia: un asistente de IA que compone mejor}
%============================================================================

\subsection{De la generación libre a la búsqueda tipada}

La construcción de un flujo de trabajo por un asistente de lenguaje es, sin un
modelo, un ejercicio de generación libre: el asistente lee las descripciones en
lenguaje natural de los nodos, adivina cuáles encajan y produce un grafo cuya
corrección solo se comprueba al ejecutarlo. Los errores típicos ---conectar la
salida de un nodo con una entrada que no puede consumirla, encadenar dos expansores
que hacen estallar el lote, dejar un nodo sin fuente que lo alimente--- son
exactamente los que el modelo de tres coordenadas caracteriza.

Expuestas al asistente ---como anotaciones del catálogo que él lee, sin cambio alguno
en los nodos---, las coordenadas convierten esa generación libre en una
\emph{búsqueda tipada}. Construir un flujo que lleve del mundo $X$ al mundo $Y$ pasa
a ser hallar un \emph{camino de flechas componibles} en el mapa de mundos: una
búsqueda en un grafo, donde cada paso respeta $\operatorname{dst}=\operatorname{src}$
y donde la cardinalidad acumulada advierte de las explosiones y de las conexiones
degeneradas. La tarea deja de ser ``inventar'' un grafo y pasa a ser ``encontrar''
uno correcto por composición.

\subsection{Por qué basta un modelo pequeño}

Este cambio de régimen tiene una consecuencia notable sobre el tamaño de modelo
necesario. La generación libre exige que el modelo lleve dentro, de forma implícita,
todo el conocimiento de cómo encajan los nodos, y que lo aplique sin equivocarse; es
una carga que solo los modelos grandes soportan con soltura. La búsqueda tipada
externaliza ese conocimiento al catálogo: las coordenadas son una \emph{gramática
composicional} explícita, y al modelo le queda la tarea, mucho más liviana, de
recorrer un espacio ya estructurado.

\prop{Con las tres coordenadas expuestas en el catálogo, la corrección
composicional de un flujo ---compatibilidad de mundos, ausencia de explosión
cardinal, ausencia de nodos aislados--- es verificable por reglas locales sobre cada
arista. Un asistente que respete esas reglas produce flujos componibles con
independencia de su tamaño; el modelo grande deja de ser un requisito para la
corrección y pasa a serlo solo para la ambición del diseño.}

En otras palabras: incluso un modelo pequeño puede construir un excelente flujo de
trabajo, porque las tres coordenadas le indican cuáles son las mejores conexiones
posibles y le vedan las imposibles. El modelo que empezó como una disciplina para
que \emph{nosotros} diseñáramos mejor los nodos resulta ser, sin cambios, la
disciplina que permite a una máquina \emph{ensamblarlos} mejor. Es la misma
gramática, leída desde los dos lados.

%============================================================================
\section{Conclusiones}
%============================================================================

Este documento partió de una dificultad concreta ---el diseño aislado de nodos y su
ciclo de retrabajo--- y de un intento honesto pero fallido de remediarlo: definir
las familias como clases de equivalencia de una relación de no-conexión. El fracaso
de ese intento en la transitividad no fue un callejón sin salida, sino la señal de
que el objeto estaba mal elegido. La teoría de conjuntos, que había inspirado el
intento, ofreció también la salida: fundar la equivalencia sobre el núcleo de una
función en lugar de sobre una relación postulada a mano.

Esa función es la \emph{firma} del nodo, y el nodo mismo dejó de ser un punto para
ser una \emph{flecha} entre mundos de datos. Tres coordenadas ---la firma de mundo,
que fija origen y destino; la cardinalidad; y la aridad de puertos que gobierna la
ramificación--- bastan para ubicar cualquier nodo en el ecosistema, clasifican desde
un \texttt{chart} hasta un \texttt{switch}, y recuperan como teorema la intuición
que la relación fallida quería como definición. Cada coordenada se justificó por el
colapso que resolvía y que la anterior no podía.

El valor del modelo es, ante todo, metodológico: transforma ``¿qué hace este nodo?''
en ``¿dónde va su flecha, y ese cruce tiene tráfico?'', y con ello convierte el
diseño rico en interacciones de un hallazgo tardío y empírico en una comprobación
previa. Que la misma estructura, expuesta a un asistente de lenguaje, convierta la
construcción de flujos en una búsqueda tipada al alcance de modelos pequeños, es la
confirmación de que las coordenadas capturan algo real sobre cómo se componen los
nodos ---y no meramente sobre cómo los nombramos---. El modelo no restringe el motor
ni toca el catálogo: solo nos enseña a mirarlo. Pero esa mirada es, precisamente, lo
que faltaba.

%============================================================================
%  Referencias
%============================================================================
\begin{thebibliography}{9}

\bibitem{hrbacek}
K.~Hrbacek and T.~Jech.
\newblock \emph{Introduction to Set Theory}, 3rd ed.
\newblock Marcel Dekker, New York, 1999.
\newblock (Relaciones de equivalencia, particiones y el núcleo de una función:
cap.~2.)

\bibitem{maclane}
S.~Mac~Lane.
\newblock \emph{Categories for the Working Mathematician}, 2nd ed.
\newblock Springer, 1998.
\newblock (El nodo como morfismo y la conexión como composición de flechas.)

\bibitem{pierce}
B.~C. Pierce.
\newblock \emph{Basic Category Theory for Computer Scientists}.
\newblock MIT Press, 1991.

\bibitem{synapse}
Mauricio Daniel Klainbard.
\newblock \emph{Synapse: un modelo de ejecución con cómputo en el flujo}.
\newblock Documento técnico, \the\year{}.
\newblock (Modelo de streaming, cardinalidad de consumo y rompedores de tubería.)

\end{thebibliography}

\end{document}
